\section{Conclusion}
\label{sec:conclusion}

\textsc{Sentinel} shows that a complete detection-to-response pipeline for a
home network, including a language model, runs on a four-core, 8\,GB-class host of the kind it targets,
and that everything except the language model and one external lookup runs in
well under a millisecond per event. Its evaluation shows how much a system can
look right while being wrong in specific, consequential ways: a detector that
misses two whole attack classes behind a respectable aggregate, false alarms
concentrated on one unseen device, a live feature extractor that fed the model
inputs it was never trained on, and a local LLM whose ATT\&CK labels are almost
never correct while its severities still drive blocking.

None of these problems was visible from the headline numbers, and each was
found by a check that is cheap to run: report per class and per device, hold
out whole devices, test live features against invariants of the training data,
pre-register the LLM test, and time every stage of the path rather than the
model alone. We offer these checks, and the harness that implements them, as
the more lasting contribution. Immediate next steps are to address the
BASHLITE flood blind spot with features that capture it, to calibrate
thresholds per device class, to make the reputation lookup asynchronous and
selective, to replace free-form LLM labelling with a validated or constrained
classifier, and to repeat the timing on physical Raspberry~Pi~5 hardware.
